\section{Conclusion}\label{sec:conclusion}

On one explicit four-dimensional $\SUtwo$ block, the gauge quotient, Haar
measure, retained--fiber dimensions, curvature coercivity, nonlinear
heat-kernel action, conditional law, and generated effective Hessian can be
placed in a single rigorous chain.  The chain is certified by exact arithmetic
and analytic inequalities: $H\succeq1/25$, the nonlinear scaled Hessian error
is below $1/400$, the fiber curvature is at least $3/80$, the conditional
transport coefficient is at most $1/3$, and one exact marginalization changes
the whitened retained Hessian by at most $1/15$.

This answers the finite question posed in the introduction.  It also sharpens
the next question.  Progress toward an iterated Yang--Mills construction no
longer depends on the algebra of this one local marginalization; it depends on
a gauge-invariant, volume-uniform control of the interacting chart complement
and on closure of the generated quasilocal action class.  Until those bridges
are proved, the result remains a restricted finite-block theorem and has no
claim to a continuum or Hamiltonian mass gap.
